% Part: sets-functions-relations
% Chapter: sets
% Section: subsets

\documentclass[../../../include/open-logic-section]{subfiles}

\begin{document}

\olfileid{sfr}{set}{sub}
\olsection{Deelverzamelingen en machtsverzamelingen}

\begin{explain}
We zullen verzamelingen vaak met elkaar willen vergelijken. Een voor de
hand liggende vergelijking is: \emph{alles wat in de ene verzameling
zit, zit ook in de andere}. Deze situatie is belangrijk genoeg om er
nieuwe notatie voor in te voeren.
\end{explain}

\begin{defn}[Deelverzameling]
Als elk !!{element} van een verzameling $A$ ook !!a{element} van~$B$
is, noemen we $A$ een \emph{deelverzameling} van~$B$ en schrijven we
$A \subseteq B$. Als $A$ geen deelverzameling van~$B$ is, schrijven
we $A \not\subseteq B$. Als $A \subseteq B$ maar $A \neq B$, schrijven
we $A \subsetneq B$ en noemen we $A$ een \emph{echte
deelverzameling} van $B$.
\end{defn}

\begin{ex}
Elke verzameling is een deelverzameling van zichzelf, en $\emptyset$
is een deelverzameling van elke verzameling. De verzameling van de even
natuurlijke getallen is een deelverzameling van de verzameling van de
natuurlijke getallen. Ook geldt $\{ a, b \} \subseteq \{ a, b, c \}$.
Daarentegen is $\{ a, b, e \}$ geen deelverzameling van
$\{ a, b, c \}$.
\end{ex}

\begin{ex}
Het getal $2$ is een !!{element} van de verzameling van de gehele
getallen, terwijl de verzameling van de even getallen een
deelverzameling is van de verzameling van de gehele getallen. Een
verzameling kan echter \emph{zowel} !!a{element} als een deelverzameling van
een andere verzameling zijn. Zo geldt $\{0\} \in \{0, \{0\}\}$ en ook
$\{0\} \subseteq \{0, \{0\}\}$.
\end{ex}

Extensionaliteit geeft een gelijkheidscriterium voor verzamelingen:
$A = B$ dan en slechts dan als elk !!{element} van~$A$ ook
!!a{element} van~$B$ is, en omgekeerd. Volgens de definitie van
‘deelverzameling’ drukt $A \subseteq B$ precies de eerste helft van dit
criterium uit: elk !!{element} van~$A$ is ook !!a{element} van~$B$.
De definitie geldt uiteraard ook met $A$ en $B$ verwisseld: $B
\subseteq A$ dan en slechts dan als elk !!{element} van~$B$ ook
!!a{element} van~$A$ is. Dat is precies wat in extensionaliteit met ‘en
omgekeerd’ wordt bedoeld. Extensionaliteit impliceert dus dat
verzamelingen gelijk zijn dan en slechts dan als ze deelverzamelingen
van elkaar zijn.

\begin{prop}
$A = B$ dan en slechts dan als zowel $A \subseteq B$ als
$B \subseteq A$.
\end{prop}

Dit is ook een geschikt moment om nog enkele handige notaties in te
voeren. In de definitie van $A$ als deelverzameling van~$B$ schreven
we ‘elk !!{element} van~$A$ is \dots,’ en vulden we de ‘$\dots$’ aan
met ‘!!a{element} van $B$’. Uitdrukkingen met deze \emph{vorm} komen
zo vaak voor dat het nuttig is er formele notatie voor in te voeren.

\begin{defn}\ollabel{forallxina}
$(\forall x \in A)\phi$ is een afkorting voor $\forall x(x \in A \lif
\phi)$. Evenzo is $(\exists x \in A)\phi$ een afkorting voor $\exists
x(x \in A \land \phi)$.
\end{defn}

Met deze notatie kunnen we zeggen dat $A \subseteq B$ dan en slechts
dan als $(\forall x \in A)x \in B$.

We gaan nu over tot een bepaald soort verzameling: de verzameling van
alle deelverzamelingen van een gegeven verzameling.

\begin{defn}[Machtsverzameling]
De verzameling die uit alle deelverzamelingen van een verzameling~$A$
bestaat, heet de \emph{machtsverzameling van}~$A$ en wordt genoteerd
als $\Pow{A}$.
  \[
    \Pow{A} = \Setabs{B}{B \subseteq A} 
  \]
\end{defn}

\begin{ex}
Welke deelverzamelingen van $\{ a, b, c \}$ zijn er? Dit zijn:
$\emptyset$, $\{a \}$, $\{b\}$, $\{c\}$, $\{a, b\}$, $\{a, c\}$,
$\{b, c\}$ en $\{a, b, c\}$. De verzameling van al deze
deelverzamelingen is $\Pow{\{a,b,c\}}$:
\[
\Pow{\{ a, b, c \}} = \{\emptyset, \{a \}, \{b\}, \{c\}, \{a, b\},
\{b, c\}, \{a, c\}, \{a, b, c\}\}
\]
\end{ex}

\begin{prob}
Geef alle deelverzamelingen van $\{a, b, c, d\}$.
\end{prob}

\begin{prob}
Toon aan dat als $A$ uit $n$ !!{element}s bestaat, $\Pow{A}$ uit
$2^n$ !!{element}s bestaat.
\end{prob}
\end{document}
